\section{符号说明}
\label{sec:symbols}

为避免将文物与其采样点混为一谈，本文分别设置文物层和采样点层索引，并按照附件中的固定顺序对14种化学成分进行编号。建模过程中使用的主要符号及其含义如表~\ref{tab:symbols}所示。

\begin{table}[H]
  \centering
  \caption{主要符号及其含义}
  \label{tab:symbols}
  \small
  \renewcommand{\arraystretch}{1.15}
  \begin{tabularx}{0.96\textwidth}{C{2.6cm} X C{3.1cm}}
    \toprule
    符号 & 含义 & 单位或取值 \\
    \midrule

    \(i,r,j\)
    & 文物、文物内部采样点及化学成分的索引
    & \(i=01,\ldots,58,\mathrm{A1},\ldots,\mathrm{A8}\)；
      \(r=1,2,\ldots\)；
      \(j=1,\ldots,14\) \\

    \(T_i\)
    & 第\(i\)件文物的玻璃类型
    & \(\mathrm{K}\)为高钾玻璃，\(\mathrm{PB}\)为铅钡玻璃 \\

    \(D_i\)
    & 第\(i\)件文物的纹饰类别
    & \(\mathrm{A},\mathrm{B},\mathrm{C}\) \\

    \(R_i\)
    & 第\(i\)件文物的颜色类别
    & 离散类别变量 \\

    \(W_i^{(A)}\)
    & 第\(i\)件文物整体表面的风化状态
    & 0为无风化，1为风化 \\

    \(W_{ir}^{(P)}\)
    & 第\(i\)件文物第\(r\)个采样点的风化等级
    & 0为未风化，1为一般风化，2为严重风化 \\

    \(x_{irj}\)
    & 第\(i\)件文物第\(r\)个采样点中第\(j\)种化学成分的原始检测比例
    & \(\%\) \\

    \(\mathbf{x}_{ir}\)
    & 第\(i\)件文物第\(r\)个采样点的14维原始成分向量，
      即\(\mathbf{x}_{ir}=(x_{ir1},\ldots,x_{ir,14})\)
    & \(\%\) \\

    \(S_{ir}\)
    & 第\(i\)件文物第\(r\)个采样点中14种化学成分比例之和，
      \(S_{ir}=\sum_{j=1}^{14}x_{irj}\)
    & \(\%\) \\

    \(V_{ir}\)
    & 采样点成分数据的有效性指示变量；
      当\(85\leq S_{ir}\leq105\)时取1，否则取0
    & \(\{0,1\}\) \\

    \(c_{irj}\)
    & 经闭合归一化后第\(j\)种化学成分的比例，
      \(c_{irj}=100x_{irj}/S_{ir}\)
    & \(\%\) \\

    \(z_{irj}\)
    & 经零值替代和中心对数比变换后第\(j\)种化学成分的变量
    & 无量纲 \\

    \(\widehat{x}_{irj}^{(0)}\)
    & 根据风化点检测数据预测得到的风化前第\(j\)种化学成分比例
    & \(\%\) \\

    \(\widehat{T}_a\)
    & 表单3中未知文物\(a\)的预测玻璃类型
    & \(\mathrm{K}\)或\(\mathrm{PB}\) \\

    \(\widehat{p}_a\)
    & 未知文物\(a\)属于铅钡玻璃的预测概率
    & \([0,1]\) \\

    \(C_{ir}^{(T)}\)
    & 在玻璃类型\(T\)内部，第\(i\)件文物第\(r\)个采样点的亚类标签
    & \(1,\ldots,K_T\) \\

    \(Q_a\)
    & 未知文物\(a\)在扰动或重采样条件下的分类稳定率
    & \([0,1]\) \\

    \(\rho_{jk}^{(t)}\)
    & 玻璃类型\(t\)中第\(j\)种与第\(k\)种化学成分之间的相关系数
    & \([-1,1]\) \\

    \(\Delta\rho_{jk}\)
    & 高钾玻璃与铅钡玻璃中成分\(j,k\)相关系数的差异
    & \([-2,2]\) \\

    \bottomrule
  \end{tabularx}
\end{table}

其中，上标\(A\)和\(P\)分别表示文物层（artifact）和采样点层（point），
不表示纹饰类别或概率。对于表面风化但存在局部未风化区域的文物，
可能出现\(W_i^{(A)}=1\)而\(W_{ir}^{(P)}=0\)的情况。
此外，\(C_{ir}^{(T)}\)中的上标\(T\)表示所属玻璃大类，并非矩阵转置。
附件中的空白项表示未检测到相应成分；计算成分总和时按0计，
进行对数比变换前则需另行实施零值替代。